\section*{Bab \ref{ch:b1:linmaps} --- Pemetaan Linear}

\begin{solution}{exo:b1:linmaps:1}
(1) \omterm{def:b1:linmaps:def}{Linear}: karena \omterm{prop:b1:vspaces:coordinates}{koordinatnya} berupa ungkapan \omterm{def:b1:linmaps:def}{linear}. (2) Tak \omterm{def:b1:linmaps:def}{linear}:
karena $u(2(1,1)) = 4 \neq 2 = 2u(1,1)$. (3) \omterm{def:b1:linmaps:def}{Linear}: karena pendiferensialan dan
perkalian dengan $X$ demikian, dan jumlah \omterm{def:b1:linmaps:def}{pemetaan linear} pun demikian. (4) \omterm{def:b1:linmaps:def}{Linear}:
karena penilaian menghormati operasi titik demi titiknya.
\end{solution}

\begin{solution}{exo:b1:linmaps:2}
\omterm{def:b1:linmaps:kerim}{Kernelnya}: pecahkanlah $x + y - z = 0$, $2x + y + z = 0$, $3x + 2y = 0$.
Dari yang ketiga, $y = -\frac{3x}{2}$; lalu yang pertama memberikan $z = x + y =
-\frac x2$; lalu periksalah pada yang kedua: $2x - \frac{3x}{2} - \frac x2 = 0$:
jadi terpenuhi. Sehingga $\ker u = \operatorname{Vect}\bigl((2, -3, -1)\bigr)$
(dengan mengambil $x = 2$), berdimensi $1$.

Lalu rank--nulitas: $\operatorname{rk} u = 3 - 1 = 2$. Adapun \omterm{def:b1:linmaps:kerim}{petanya}: direntang oleh
\omterm{def:b1:linmaps:kerim}{peta} \omterm{def:b1:vspaces:free}{basis} kanoniknya, $u(e_1) = (1,2,3)$, $u(e_2) =
(1,1,2)$, $u(e_3) = (-1,1,0)$; dengan dua yang pertama \omterm{def:b1:vspaces:free}{bebas}, dan \omterm{def:b1:linmaps:rank}{ranknya}
$2$: sehingga basisnya $\bigl((1,2,3), (1,1,2)\bigr)$.

Jadi tak \omterm{def:b1:logic:inj}{injektif} (karena $\ker \neq \{0\}$), dan tak \omterm{def:b1:logic:inj}{surjektif} (karena \omterm{def:b1:linmaps:rank}{ranknya} $2 < 3$):
yang selaras dengan \cref{cor:b1:linmaps:samedim}.
\end{solution}

\begin{solution}{exo:b1:linmaps:3}
\emph{\omterm{def:b1:linmaps:kerim}{Kernelnya}:} $P = P'$ memaksa $\deg P = \deg P'$ kecuali bila $P = 0$;
padahal $\deg P' < \deg P$ untuk $P \neq 0$: sehingga $\ker u = \{0\}$, dan $u$,
sebagai endomorfisma \omterm{def:b1:logic:inj}{injektif} pada $\R_n[X]$ yang \omterm{def:b1:findim:def}{berdimensi hingga}, merupakan
isomorfisma (\cref{cor:b1:linmaps:samedim}).

\emph{Inversnya:} misalkan $v(P) = P + P' + P'' + \dots + P^{(n)}$ (yaitu
jumlah yang hingga pada $\R_n[X]$). Maka
\[
v\bigl(u(P)\bigr) = \sum_{k=0}^{n} (P - P')^{(k)}
= \sum_{k=0}^{n} P^{(k)} - \sum_{k=0}^{n} P^{(k+1)}
= P - P^{(n+1)} = P ,
\]
secara teleskopis, karena $P^{(n+1)} = 0$. Jadi $v = u^{-1}$.
\end{solution}

\begin{solution}{exo:b1:linmaps:4}
Di sini $\operatorname{im}(v \circ u) = v(\operatorname{im} u) \subseteq
\operatorname{im} v$: sehingga \omterm{def:b1:linmaps:rank}{ranknya} $\leq \operatorname{rk} v$. Dan $v$
yang dibatasi pada $\operatorname{im} u$ berpeta
$\operatorname{im}(vu)$, sehingga dengan rank--nulitas di dalam
$\operatorname{im} u$: $\operatorname{rk}(vu) \leq
\dim\operatorname{im} u = \operatorname{rk} u$.
\end{solution}

\begin{solution}{exo:b1:linmaps:5}
Persamaan $u^2 = 0$ berarti $u(u(x)) = 0$ untuk setiap $x$: sehingga setiap $u(x)$ terletak di
$\ker u$, yakni $\operatorname{im} u \subseteq \ker u$. Lalu menurut
rank--nulitas:
\[
\dim E = \dim\ker u + \operatorname{rk} u \geq 2\operatorname{rk} u .
\]
Adapun contoh kesamaannya di $\R^2$: $u(x, y) = (y, 0)$: karena $u^2 = 0$,
$\operatorname{rk} u = 1 = \frac{\dim E}{2}$.
\end{solution}

\begin{solution}{exo:b1:linmaps:6}
Di sini $(pq)^2 = pqpq = ppqq = pq$ (menurut komutasinya): jadi sebuah \omterm{def:b1:linmaps:projection}{proyeksi}
(\cref{thm:b1:linmaps:projchar}).

\omterm{def:b1:linmaps:kerim}{Petanya}: $\operatorname{im}(pq) \subseteq \operatorname{im} p$, dan
$= \operatorname{im}(qp) \subseteq \operatorname{im} q$: sehingga termuat
di irisannya. Sebaliknya, jika $x \in \operatorname{im} p \cap
\operatorname{im} q$, maka $p(x) = x$ dan $q(x) = x$ (karena titik tetapnya
mencirikan \omterm{def:b1:linmaps:kerim}{peta} sebuah \omterm{def:b1:linmaps:projection}{proyeksi}), sehingga $pq(x) = x$: jadi $x \in
\operatorname{im}(pq)$.

\omterm{def:b1:linmaps:kerim}{Kernelnya}: $\ker p \subseteq \ker(qp) = \ker(pq)$ dan demikian pula $\ker q
\subseteq \ker(pq)$: sehingga jumlahnya termuat. Sebaliknya, misalkan $pq(x) =
0$, lalu tulislah
\[
x = \underbrace{q(x)}_{\in\, \ker p} +
\underbrace{(x - q(x))}_{\in\, \ker q} :
\]
maka suku pertamanya memenuhi $p(q(x)) = 0$, sehingga ia terletak di $\ker p$; sedangkan
yang kedua berada di $\ker q$ karena $q(x - q(x)) = q(x) - q^2(x) = 0$. Jadi
$x \in \ker p + \ker q$.
\end{solution}

\begin{solution}{exo:b1:linmaps:7}
Perhatikanlah lebih dulu inklusi umumnya $\ker u \subseteq \ker u^2$ dan
$\operatorname{im} u^2 \subseteq \operatorname{im} u$, lalu menurut
rank--nulitas, (2) $\iff$ (3) (karena \omterm{def:b1:linmaps:kerim}{kernel} yang sama $\iff$ \omterm{def:b1:linmaps:rank}{rank} yang sama
$\iff$ \omterm{def:b1:linmaps:kerim}{peta} yang sama, mengingat inklusinya).

(1 $\Rightarrow$ 2): misalkan $u^2(x) = 0$; maka $u(x) \in \ker u \cap
\operatorname{im} u = \{0\}$, sehingga $x \in \ker u$.

(2 $\Rightarrow$ 1): menurut Grassmann dan rank--nulitas, $\dim(\ker u +
\operatorname{im} u) = \dim\ker u + \operatorname{rk} u -
\dim(\ker u \cap \operatorname{im} u) = \dim E - \dim(\ker u \cap
\operatorname{im} u)$: sehingga jumlahnya adalah $E$ jika dan hanya jika irisannya
$\{0\}$. Misalkan $y \in \ker u \cap \operatorname{im} u$: maka $y = u(x)$ dan
$u(y) = 0$, sehingga $u^2(x) = 0$, jadi (menurut (2)) $u(x) = 0$: sehingga $y = 0$. Jadi
$E = \ker u \oplus \operatorname{im} u$.
\end{solution}

\begin{solution}{exo:b1:linmaps:8}
Jika $\varphi = 0$: maka $\ker\varphi = E \subseteq \ker\psi$ memaksa
$\psi = 0 = 0\cdot\varphi$. Sedangkan jika $\varphi \neq 0$: maka $\ker\varphi$ sebuah
\omterm{def:b1:linmaps:forms}{hiperbidang}; pilihlah $a \notin \ker\varphi$. Tetapkanlah $\lambda =
\frac{\psi(a)}{\varphi(a)}$. Maka bentuk $\psi - \lambda\varphi$
lenyap pada $\ker\varphi$ (karena keduanya demikian, menurut inklusinya) dan di $a$:
sehingga ia lenyap pada $\ker\varphi \oplus Ka = E$. Jadi $\psi =
\lambda\varphi$.
\end{solution}

\begin{solution}{exo:b1:linmaps:9}
Andaikan $\lambda_0 x + \lambda_1 u(x) + \dots + \lambda_{n-1}
u^{n-1}(x) = 0$. Terapkanlah $u^{n-1}$: maka semua suku berfaktor $u^{\geq
n}$ mati, sehingga menyisakan $\lambda_0 u^{n-1}(x) = 0$, jadi $\lambda_0 = 0$.
Lalu terapkanlah $u^{n-2}$ pada relasi yang tersisa: $\lambda_1 u^{n-1}(x) =
0$, sehingga $\lambda_1 = 0$; dan seterusnya. Jadi keluarganya \omterm{def:b1:vspaces:free}{bebas}; dan karena berukuran
$n = \dim E$, ia sebuah \omterm{def:b1:vspaces:free}{basis} (\cref{prop:b1:findim:twoofthree}). (Pada
\omterm{def:b1:vspaces:free}{basis} ini, $u$ bertindak sebagai geseran --- yaitu model kenilpotenan
yang maksimal.)
\end{solution}

\begin{solution}{exo:b1:linmaps:10}
\begin{enumerate}
  \item \omterm{def:b1:logic:map}{Pemetaan} $f \circ f = -\mathrm{id}$ bersifat \omterm{def:b1:logic:inj}{bijektif}, sehingga $f$ demikian
        (\cref{prop:b1:logic:comp} yang disesuaikan: karena $f$ mempunyai invers
        dua sisinya $-f$). Jika $f(x) = \lambda x$ dengan $x \neq 0$:
        maka menerapkan $f$, $-x = \lambda^2 x$, sehingga $\lambda^2 = -1$:
        yang mustahil di $\R$.
  \item Bangunlah keluarganya secara rakus. Ambillah $x_1 \neq 0$: maka $(x_1,
        f(x_1))$ \omterm{def:b1:vspaces:free}{bebas} menurut (1). Jika $\operatorname{Vect}$ atas
        keluarga yang berjalan $\bigl(x_1, f(x_1), \dots, x_k,
        f(x_k)\bigr)$, yang kita sebut $V_k$ --- yaitu \omterm{def:b1:vspaces:subspace}{subruang} yang stabil di bawah $f$
        (karena setiap pembangunnya terpetakan ke pembangun lain atau lawannya:
        $f(f(x_i)) = -x_i$) --- bukan seluruh $E$, maka pilihlah $x_{k+1}
        \notin V_k$. \emph{Klaimnya: keluarga yang diperbesar itu \omterm{def:b1:vspaces:free}{bebas}.}
        Andaikan $\alpha x_{k+1} + \beta f(x_{k+1}) + v = 0$ dengan $v
        \in V_k$ dan $(\alpha, \beta) \neq (0,0)$. Terapkanlah $f$:
        $\alpha f(x_{k+1}) - \beta x_{k+1} + f(v) = 0$ dengan $f(v)
        \in V_k$. Lalu hapuslah $f(x_{k+1})$ di antara kedua relasinya
        (kalikanlah yang pertama dengan $\alpha$, yang kedua dengan $-\beta$, lalu
        jumlahkan):
        \[
        (\alpha^2 + \beta^2)\, x_{k+1} \in V_k ,
        \]
        dan $\alpha^2 + \beta^2 \neq 0$ memaksa $x_{k+1} \in V_k$:
        yang bertentangan. Jadi konstruksinya berlanjut, sambil menambahkan vektor
        \emph{dua sekaligus}, sampai $V_k = E$: sehingga keluarga akhirnya sebuah
        \omterm{def:b1:vspaces:free}{basis} berukuran genap, jadi $n$ genap.
\end{enumerate}
\end{solution}

\begin{solution}{exo:b1:linmaps:11}
Batas atasnya: untuk setiap $x$, $(u + v)(x) = u(x) + v(x) \in
\operatorname{im} u + \operatorname{im} v$, sehingga
\[
\operatorname{rk}(u + v)
\leq \dim(\operatorname{im} u + \operatorname{im} v)
\leq \operatorname{rk} u + \operatorname{rk} v
\]
(menurut Grassmann, \cref{thm:b1:findim:grassmann}). Adapun batas bawahnya: terapkanlah
batas atasnya pada pasangan $(u + v, -v)$, yang jumlahnya $u$:
\[
\operatorname{rk} u \leq \operatorname{rk}(u + v) +
\operatorname{rk}(-v) = \operatorname{rk}(u + v) +
\operatorname{rk} v,
\]
sehingga $\operatorname{rk} u - \operatorname{rk} v \leq
\operatorname{rk}(u+v)$; lalu menukar $u$ dan $v$ memberikan nilai
mutlaknya.
\end{solution}

\begin{solution}{exo:b1:linmaps:12}
\emph{Rumus yang persis.} Misalkan $v'$ pembatasan $v$ pada
\omterm{def:b1:vspaces:subspace}{subruang} $\operatorname{im} w$. Maka \omterm{def:b1:linmaps:kerim}{petanya} adalah $v(w(F)) =
\operatorname{im}(v \circ w)$, dan \omterm{def:b1:linmaps:kerim}{kernelnya} adalah $\ker v \cap
\operatorname{im} w$. Lalu rank--nulitas bagi $v'$ pada ruang
$\operatorname{im} w$:
\[
\operatorname{rk} w = \dim\operatorname{im} w
= \operatorname{rk}(v \circ w) + \dim(\ker v \cap
\operatorname{im} w) .
\]

\emph{Frobenius.} Terapkanlah rumus yang persis itu dua kali, pada $w$ dan pada
$w \circ u$:
\[
\operatorname{rk} w - \operatorname{rk}(vw)
= \dim\bigl(\ker v \cap \operatorname{im} w\bigr),
\qquad
\operatorname{rk}(wu) - \operatorname{rk}(vwu)
= \dim\bigl(\ker v \cap \operatorname{im}(wu)\bigr) .
\]
Dan karena $\operatorname{im}(w \circ u) \subseteq \operatorname{im}
w$, irisan yang kedua termuat di yang pertama, sehingga
dimensinya tak lebih besar:
\[
\operatorname{rk}(wu) - \operatorname{rk}(vwu)
\;\leq\; \operatorname{rk} w - \operatorname{rk}(vw) ,
\]
yang tersusun ulang menjadi ketaksamaan Frobenius. Adapun dengan $w =
\mathrm{id}_F$ (yang \omterm{def:b1:linmaps:rank}{ber-rank} $\dim F$, dan $\operatorname{im}\,
\mathrm{id}_F = F$): $\operatorname{rk} v + \operatorname{rk} u
\leq \dim F + \operatorname{rk}(vu)$, yaitu ketaksamaan Sylvester ---
yang dibuktikan lagi, dalam bentuk matriks, pada \cref{exo:b1:matrices:10}.
\end{solution}

\begin{solution}{pb:b1:linmaps:1}
\textbf{1.} Di sini $(\mathrm{id} - p)^2 = \mathrm{id} - 2p + p^2 =
\mathrm{id} - p$: jadi sebuah proyektor. Jika $y = x - p(x)$, maka $p(y) =
p(x) - p^2(x) = 0$, dan sebaliknya $x \in \ker p$ memberikan $x =
(\mathrm{id} - p)(x)$: sehingga $\operatorname{im}(\mathrm{id} - p) = \ker
p$. Dan $(\mathrm{id} - p)(x) = 0 \iff p(x) = x \iff x \in
\operatorname{im} p$ (menurut titik tetapnya,
\cref{thm:b1:linmaps:projchar}): sehingga $\ker(\mathrm{id} - p) =
\operatorname{im} p$.

\textbf{2.} Di sini $(\lambda\,\mathrm{id} + \mu p)^2 =
\lambda^2\,\mathrm{id} + (2\lambda\mu + \mu^2)\,p$. Adapun pasangan
$(\mathrm{id}, p)$ \omterm{def:b1:vspaces:free}{bebas} di $\mathcal{L}(E)$: karena $p = c\,
\mathrm{id}$ akan memberikan $c^2 = c$, sehingga $p = 0$ atau $\mathrm{id}$,
yang tersingkirkan. Lalu dengan mengenali koefisiennya, \omterm{def:b1:logic:map}{pemetaannya} merupakan proyektor jika dan hanya jika
$\lambda^2 = \lambda$ dan $2\lambda\mu + \mu^2 = \mu$. Untuk
$\lambda = 0$: $\mu \in \{0, 1\}$. Untuk $\lambda = 1$: $\mu^2 +
\mu = 0$, sehingga $\mu \in \{0, -1\}$. Jadi tepat empat proyektor pada
bidangnya: $0$, $p$, $\mathrm{id}$, dan $\mathrm{id} - p$.

\textbf{3.} Di sini $(\lambda\,\mathrm{id} + \mu p)(\lambda'\,\mathrm{id}
+ \mu' p) = \lambda\lambda'\,\mathrm{id} + (\lambda\mu' +
\mu\lambda' + \mu\mu')\,p$: sehingga bidangnya stabil terhadap
komposisi. Dan karena $p^k = p$ untuk setiap $k \geq 1$, maka untuk $Q =
\sum_k a_k X^k$:
\[
Q(p) = a_0\,\mathrm{id} + \Bigl(\sum_{k \geq 1} a_k\Bigr) p
= Q(0)\,\mathrm{id} + \bigl(Q(1) - Q(0)\bigr)\,p .
\]

\textbf{4.} Pada $\operatorname{im} p$ (yang di situ $p$ bertindak sebagai
identitas), $\lambda\,\mathrm{id} + \mu p$ mengalikan dengan $\lambda
+ \mu$; sedangkan pada $\ker p$, dengan $\lambda$. Dan karena $E = \operatorname{im} p
\oplus \ker p$, \omterm{def:b1:logic:map}{pemetaannya} bersifat \omterm{def:b1:logic:inj}{bijektif} jika dan hanya jika $\lambda \neq 0$ dan
$\lambda + \mu \neq 0$. Lalu memecahkan $\lambda\alpha = 1$,
$\lambda\beta + \mu\alpha + \mu\beta = 0$ pada aturan komposisi
pertanyaan 3:
\[
(\lambda\,\mathrm{id} + \mu p)^{-1}
= \frac1\lambda\,\mathrm{id} -
\frac{\mu}{\lambda(\lambda + \mu)}\,p ,
\]
yang tindakannya berupa $1/\lambda$ pada $\ker p$ dan $1/(\lambda +
\mu)$ pada $\operatorname{im} p$, sebagaimana seharusnya.

\textbf{5.} Tulislah $F = \operatorname{im} p = \operatorname{im}
p'$. Maka untuk setiap $x$, $p'(x) \in F$ dan $p$ menetapkan $F$ titik demi titik:
sehingga $p(p'(x)) = p'(x)$, yakni $p\,p' = p'$; dan secara simetris $p'\,p =
p$. Jadi ketika dua \omterm{def:b1:linmaps:projection}{proyeksi} berbagi \omterm{def:b1:linmaps:kerim}{petanya}, yang diterapkan
\emph{lebih dulu} yang memutuskan: karena keluarannya sudah terletak di $F$, yang di situ
\omterm{def:b1:linmaps:projection}{proyeksi} yang luar bertindak sebagai identitas dan tak mengubah apa pun.

\textbf{6.} Di sini $(p + q)^2 = p^2 + pq + qp + q^2 = (p + q) + pq +
qp$, sehingga $p + q$ yang proyektor memaksa $pq + qp = 0$. Lalu susunlah di
kiri dengan $p$: $pq + pqp = 0$; dan di kanan dengan $p$: $pqp + qp
= 0$. Lalu menguranginya, $pq = qp$; sehingga $pq + qp = 2pq = 0$ dan karena
karakteristiknya bukan $2$: maka $pq = qp = 0$.

\textbf{7.} Dengan $pq = qp = 0$, penjabaran yang sama memberikan $(p +
q)^2 = p + q$. Adapun \omterm{def:b1:linmaps:kerim}{petanya}: $\operatorname{im}(p + q) \subseteq
\operatorname{im} p + \operatorname{im} q$ selalu. Sebaliknya,
untuk $x \in \operatorname{im} p$: $q(x) = q(p(x)) = 0$, sehingga $(p +
q)(x) = p(x) = x$ dan $x \in \operatorname{im}(p+q)$; demikian pula bagi
$\operatorname{im} q$. Adapun kelangsungannya: $x \in \operatorname{im} p
\cap \operatorname{im} q$ memberikan $x = p(x) = p(q(x)) = 0$.
Adapun \omterm{def:b1:linmaps:kerim}{kernelnya}: jika $p(x) + q(x) = 0$, maka menerapkan $p$ memberikan $p(x) +
p(q(x)) = p(x) = 0$, dan menerapkan $q$ memberikan $q(x) = 0$: sehingga $\ker(p
+ q) = \ker p \cap \ker q$ (dengan inklusi baliknya yang jelas).

\textbf{8.} Fungsi $p - q$ merupakan proyektor jika dan hanya jika $\mathrm{id} - (p - q) =
(\mathrm{id} - p) + q$ demikian (menurut pertanyaan 1 dua kali). Lalu menurut pertanyaan
6--7 yang diterapkan pada proyektor $\mathrm{id} - p$ dan $q$, ini
berlaku jika dan hanya jika $(\mathrm{id} - p)q = q(\mathrm{id} - p) = 0$, yakni
jika dan hanya jika $pq = q$ dan $qp = q$.

\textbf{9.} Persamaan $pq = q$ berarti $p$ menetapkan setiap $q(x)$, yakni
$\operatorname{im} q \subseteq \ker(p - \mathrm{id}) =
\operatorname{im} p$. Sedangkan $qp = q$ berarti $q\bigl((\mathrm{id} -
p)(x)\bigr) = 0$ untuk setiap $x$, yakni $q$ lenyap pada
$\operatorname{im}(\mathrm{id} - p) = \ker p$: sehingga $\ker p \subseteq
\ker q$. Kedua langkahnya berupa kesetaraan: jadi urutan $q \leq p$ mengatakan bahwa
$q$ memproyeksikan pada \omterm{def:b1:linmaps:kerim}{peta} yang lebih kecil, sepanjang \omterm{def:b1:linmaps:kerim}{kernel} yang lebih besar.

\textbf{10.} Dengan menjabarkannya, $(\mathrm{id} - p)(\mathrm{id} - q) =
\mathrm{id} - p - q + pq = \mathrm{id} - r$. Adapun proyektor
$\mathrm{id} - p$ dan $\mathrm{id} - q$ berkomutasi, sehingga menurut
\cref{exo:b1:linmaps:6} hasil kalinya $\mathrm{id} - r$ merupakan
proyektor pada $\operatorname{im}(\mathrm{id} - p) \cap
\operatorname{im}(\mathrm{id} - q) = \ker p \cap \ker q$ sepanjang
$\ker(\mathrm{id} - p) + \ker(\mathrm{id} - q) =
\operatorname{im} p + \operatorname{im} q$. Lalu menurut pertanyaan 1, $r =
\mathrm{id} - (\mathrm{id} - r)$ menjadi proyektor dengan
$\operatorname{im} r = \operatorname{im} p + \operatorname{im}
q$ dan $\ker r = \ker p \cap \ker q$.

\textbf{11.} \omterm{def:b1:logic:map}{Pemetaan} $p_i$ terdefinisi dengan baik (menurut ketunggalan
penguraiannya) dan \omterm{def:b1:linmaps:def}{linear} (karena penguraian $x + \lambda y$
merupakan jumlah penguraiannya, kembali menurut ketunggalannya). Untuk $x_i
\in F_i$ penguraiannya adalah $x_i$ sendiri, sehingga $p_i(x_i) = x_i$:
jadi $p_i^2 = p_i$, dan $p_j(x_i) = 0$ untuk $j \neq i$: sehingga $p_i p_j = 0$
(karena $p_j(x) \in F_j$). Lalu menjumlahkan komponennya, $\sum_i p_i =
\mathrm{id}$. Akhirnya $\operatorname{im} p_i = F_i$ dan $\ker
p_i = \bigoplus_{j \neq i} F_j$.

\textbf{12.} Di sini $p_i = p_i \circ \mathrm{id} = p_i\sum_j p_j =
p_i^2 + \sum_{j \neq i} p_i p_j = p_i^2$: sehingga setiap $p_i$ sebuah
proyektor. Lalu setiap $x = \mathrm{id}(x) = \sum_i p_i(x)$ terletak di
$\sum_i \operatorname{im} p_i$: sehingga \omterm{def:b1:linmaps:kerim}{petanya} berjumlah $E$.
Adapun kelangsungannya: andaikan $y_1 + \dots + y_k = 0$ dengan $y_i \in
\operatorname{im} p_i$, sehingga $p_i(y_i) = y_i$. Terapkanlah $p_j$: maka $p_j
(y_i) = p_j p_i (y_i) = 0$ untuk $i \neq j$, sehingga $0 = p_j\bigl(\sum
y_i\bigr) = y_j$, untuk setiap $j$. Jadi $E = \bigoplus_i
\operatorname{im} p_i$.

\textbf{13.} Di sini $q = \mathrm{id} - p$, lalu pertanyaan 1 memberikan $pq =
p - p^2 = 0 = qp$ secara langsung: jadi bagi dua proyektor, berjumlah
identitas sudah memaksa keortogonalan pasangannya.

\textbf{14.} Di sini $p + q = \mathrm{id} - r$ dengan $r$ sebuah proyektor, dan
$(\mathrm{id} - r)$ merupakan proyektor (pertanyaan 1): sehingga $p + q$ sebuah
proyektor, lalu pertanyaan 6 memberikan $pq = qp = 0$. Lalu menurut kesimetriannya (karena $q +
r = \mathrm{id} - p$ dan $p + r = \mathrm{id} - q$), semua
hasil kali berpasangannya lenyap, dan pertanyaan 12 menyimpulkannya: $E =
\operatorname{im} p \oplus \operatorname{im} q \oplus
\operatorname{im} r$, secara otomatis.

\textbf{15.} \emph{Lemanya.} Lewat induksi dengan Grassmann
(\cref{thm:b1:findim:grassmann}):
\[
\dim(F_1 + \dots + F_k) \leq \dim(F_1 + \dots + F_{k-1}) + \dim
F_k \leq \dots \leq \sum_i \dim F_i .
\]
Jika totalnya berupa kesamaan, maka setiap langkahnya demikian: $(F_1 + \dots +
F_{j-1}) \cap F_j = \{0\}$ untuk setiap $j$, sehingga sebuah relasi $y_1 +
\dots + y_k = 0$ (dengan $y_i \in F_i$) runtuh dari kanannya: $y_k
\in (F_1 + \dots + F_{k-1}) \cap F_k = \{0\}$, lalu $y_{k-1} =
0$, dan seterusnya: jadi jumlahnya langsung. Sebaliknya \omterm{def:b1:vspaces:sum}{jumlah langsung} mempunyai
dimensi yang aditif (rangkaikanlah basisnya). \emph{Penerapannya:}
$x = \sum_i p_i(x)$ menunjukkan $E = \sum_i \operatorname{im} p_i$,
sehingga $n \leq \sum_i \operatorname{rk} p_i$; lalu hipotesisnya memberikan
kesamaan, jadi kelangsungannya. Adapun hasil kalinya: tetapkanlah $j$ dan $y \in
\operatorname{im} p_j$. Maka $y = \sum_i p_i(y)$ dengan $p_i(y)
\in \operatorname{im} p_i$, sedangkan $y = y$ juga sebuah
penguraian (dengan komponen $j$ belaka); sehingga ketunggalannya memaksa $p_i(y)
= 0$ untuk $i \neq j$. Lalu diterapkan pada $y = p_j(x)$: $p_i p_j = 0$.

\textbf{16.} Jika $u^k(x) = 0$ maka $u^{k+1}(x) = u(0) = 0$. Dan
$\operatorname{im} u^{k+1} = u^k\bigl(u(E)\bigr) \subseteq
u^k(E) = \operatorname{im} u^k$.

\textbf{17.} Anggaplah $\ker u^{r} = \ker u^{r+1}$ lalu misalkan $x \in
\ker u^{r+2}$: maka $u(x) \in \ker u^{r+1} = \ker u^{r}$, sehingga
$u^{r+1}(x) = 0$: jadi $x \in \ker u^{r+1}$. Digabung dengan pertanyaan 16,
$\ker u^{r+1} = \ker u^{r+2}$, dan lewat induksi semua
\omterm{def:b1:linmaps:kerim}{kernel} berikutnya berimpit dengan $\ker u^{r}$. Adapun bagi \omterm{def:b1:linmaps:kerim}{petanya}: rank--nulitas
memberikan $\dim\operatorname{im} u^k = n - \dim\ker u^k$, sehingga
dimensi \omterm{def:b1:linmaps:kerim}{petanya} membeku persis ketika dimensi \omterm{def:b1:linmaps:kerim}{kernelnya} juga demikian,
lalu dengan inklusi pertanyaan 16, dimensi yang sama berarti
\omterm{def:b1:vspaces:subspace}{subruang} yang sama (\cref{thm:b1:findim:subspaces}).

\textbf{18.} Barisan $\bigl(\dim\ker u^k\bigr)_k$ bersifat
tidak turun dengan nilai di $\intint{0}{n}$; sehingga ia tak dapat naik
tegas $n + 1$ kali, jadi suatu $\dim\ker u^{r} = \dim\ker
u^{r+1}$ dengan $r \leq n$, sehingga $\ker u^{r} = \ker u^{r+1}$
(menurut inklusinya ditambah dimensi yang sama). Ambillah $r$ yang terkecil; lalu pertanyaan 17
membekukan segalanya mulai $r$, termasuk \omterm{def:b1:linmaps:kerim}{petanya}.

\textbf{19.} Irisannya: misalkan $x \in \ker u^{r} \cap
\operatorname{im} u^{r}$, katakanlah $x = u^{r}(y)$ dengan $u^{r}(x) =
0$. Maka $u^{2r}(y) = 0$, dan $\ker u^{2r} = \ker u^{r}$
(pertanyaan 17), sehingga $x = u^{r}(y) = 0$. Adapun dimensinya: rank--nulitas
bagi $u^{r}$ memberikan $\dim\ker u^{r} + \dim\operatorname{im} u^{r}
= n$; lalu dengan irisan yang sepele, Grassmann membuat jumlahnya menjadi
\omterm{def:b1:vspaces:subspace}{subruang} berdimensi $n$: sehingga $E = \ker u^{r} \oplus
\operatorname{im} u^{r}$.

\textbf{20.} Kestabilannya: $u^{r}(u(x)) = u(u^{r}(x)) = 0$ untuk $x
\in \ker u^{r}$; dan $u(u^{r}(y)) = u^{r}(u(y)) \in
\operatorname{im} u^{r}$. Pada $N = \ker u^{r}$: $(u|_N)^{r} = 0$
menurut definisi $N$: jadi nilpoten. Sedangkan pada $I = \operatorname{im}
u^{r}$: $\ker(u|_I) = \ker u \cap I \subseteq \ker u^{r} \cap I
= \{0\}$, sehingga $u|_I$ merupakan endomorfisma \omterm{def:b1:logic:inj}{injektif} pada
$I$ yang \omterm{def:b1:findim:def}{berdimensi hingga}, jadi \omterm{def:b1:logic:inj}{bijektif}
(\cref{cor:b1:linmaps:samedim}).

\textbf{21.} Misalkan $x = a + b$ dengan $a \in N$, $b \in I$. Maka
$u(x) = u(a) + u(b)$ dengan $u(a) \in N$ dan $u(b) \in I$
(pertanyaan 20): dan inilah \emph{persis} penguraian $u(x)$, sehingga
$\pi(u(x)) = u(a) = u(\pi(x))$: jadi $\pi u = u\pi$.

\textbf{22.} Di sini $u^2(x,y,z) = u(y, 0, z) = (0, 0, z)$ dan
$u^3(x,y,z) = u(0,0,z) = (0,0,z) = u^2(x,y,z)$. Adapun \omterm{def:b1:linmaps:kerim}{kernelnya}: $\ker u
= \{y = z = 0\} = \operatorname{Vect}(e_1)$, $\ker u^2 = \{z =
0\} = \operatorname{Vect}(e_1, e_2)$, $\ker u^3 = \ker u^2$:
sehingga stabil pada $r = 2$. Adapun \omterm{def:b1:linmaps:kerim}{petanya}: $\operatorname{im} u =
\operatorname{Vect}(e_1, e_3)$, $\operatorname{im} u^2 =
\operatorname{Vect}(e_3)$. Lalu Fitting: $\R^3 =
\operatorname{Vect}(e_1, e_2) \oplus \operatorname{Vect}(e_3)$,
dan $\pi(x, y, z) = (x, y, 0)$. Periksa: $\pi u(x,y,z) = \pi(y, 0,
z) = (y, 0, 0)$ dan $u\pi(x,y,z) = u(x, y, 0) = (y, 0, 0)$:
jadi sama. Pada faktor pertamanya $u(x, y, 0) = (y, 0, 0)$, yang
berkuadrat $0$: jadi nilpoten; sedangkan pada yang kedua $u(0,0,z) = (0,0,z)$:
yaitu identitasnya, jadi \omterm{def:b1:logic:inj}{bijektif}.

\textbf{23.} Jika $u^m = 0$ maka $\ker u^m = E$; dan karena \omterm{def:b1:linmaps:kerim}{kernelnya}
membeku mulai $r$, $\ker u^{r} = \ker u^{\max(m, r)} = E$.
Sebaliknya $\ker u^{r} = E$ berarti $u^{r} = 0$. Dan $\ker u^{r} =
E \iff$ proyektor Fittingnya memetakan pada $E$ sepanjang $\{0\}$, yakni
$\pi = \mathrm{id}$. Akhirnya $r \leq n$ (pertanyaan 18) memberikan:
bahwa setiap endomorfisma nilpoten memenuhi $u^{n} = 0$ --- sehingga
indeks kenilpotenannya tak pernah melampaui dimensinya.

\textbf{24.} Misalkan $m$ sebuah indeks kenilpotenan $u|_A$: maka $A
\subseteq \ker u^{m} \subseteq \ker u^{\max(m,r)} = \ker u^{r}$.
Dan karena $u|_B$ \omterm{def:b1:logic:inj}{bijektif}, $B = u(B) = u^{k}(B) \subseteq
\operatorname{im} u^{k}$ untuk setiap $k$, khususnya $B
\subseteq \operatorname{im} u^{r}$. Lalu
\[
n = \dim A + \dim B \leq \dim\ker u^{r} +
\dim\operatorname{im} u^{r} = n :
\]
jadi kedua inklusinya berupa kesamaan dimensi, sehingga kesamaan
\omterm{def:b1:vspaces:subspace}{subruang}: $A = \ker u^{r}$, $B = \operatorname{im} u^{r}$.

\textbf{25.} (i) Bagian III merupakan sebuah kamus: bahwa pembelahan $E = F_1
\oplus \dots \oplus F_k$ tepat bersesuaian dengan keluarga
proyektor dengan $\sum p_i = \mathrm{id}$ dan $p_i p_j = 0$, dengan
$F_i$ sebagai \omterm{def:b1:linmaps:kerim}{petanya}. (ii) Untuk $k = 3$ \omterm{def:b1:vspaces:sum}{pelengkapnya}
$\mathrm{id} - p_i$ sendirinya berupa proyektor, yang menutup
argumennya tanpa hipotesis tambahan; sedangkan untuk $k$ yang umum kita memerlukan
$\sum_i \operatorname{rk} p_i \leq n$, yaitu ketaksamaan yang diberikan
trace Tahun~2 secara cuma-cuma (karena $\operatorname{tr} p = \operatorname{rk}
p$ bagi sebuah proyektor, dan tracenya berjumlah $\operatorname{tr}
\mathrm{id} = n$). (iii) Adapun \cref{exo:b1:linmaps:7} merupakan lema
Fitting pada kasus yang sudah stabil $r \leq 1$; sedangkan secara umum kita
membiarkan rantai \omterm{def:b1:linmaps:kerim}{kernel} dan \omterm{def:b1:linmaps:kerim}{petanya} membeku, yang memakan paling banyak $n$
langkah. (iv) Pada jilid Tahun~2, bila diterapkan pada $u - \lambda\,
\mathrm{id}$, faktor nilpotennya menjadi ruang eigen yang tersamaratakan
di $\lambda$ dan proyektor Bagian III menjadi
proyektor spektral teori reduksinya. Adapun teorema
Bagian IV adalah \emph{lema Fitting}.

\end{solution}
